% Created 2019-01-21 Mon 12:31
\documentclass[11pt]{article}
\usepackage[utf8]{inputenc}
\usepackage[T1]{fontenc}
\usepackage{fixltx2e}
\usepackage{graphicx}
\usepackage{longtable}
\usepackage{float}
\usepackage{wrapfig}
\usepackage{rotating}
\usepackage[normalem]{ulem}
\usepackage{amsmath}
\usepackage{textcomp}
\usepackage{marvosym}
\usepackage{wasysym}
\usepackage{amssymb}
\usepackage{hyperref}
\tolerance=1000
\usepackage{csquotes}
\usepackage{graphicx}
\usepackage{amsmath}
\usepackage{mathtools}
\usepackage{algpseudocode}
\usepackage[margin=0.75in]{geometry}
\usepackage{nth}
\usepackage{tikz}
\usetikzlibrary{positioning,chains,fit,shapes,shapes.misc,calc,arrows,patterns,decorations.pathreplacing}
\usepackage{ccicons}
\usepackage{calc}
\usepackage{fancyhdr}
\usepackage{hyperref}
\usepackage{nopageno}
\usepackage{enumitem}
\usepackage{multicol}
\usepackage{braket}
\definecolor{ans}{rgb}{0,.5,0}%{0.545,0.27,0.074}
\def\ans#1{{\color{ans}#1}}
\algnewcommand\algorithmicforeach{\textbf{for each}}
\algdef{S}[FOR]{ForEach}[1]{\algorithmicforeach\ #1\ \algorithmicdo}
\setlength{\columnseprule}{1pt}
\newcommand{\bonusmarksamt}[1]{\textbf{[#1~BONUS~marks]}}
\newcommand{\bonusmarkamt}[0]{\textbf{[1~BONUS~mark]}}
\newcommand{\marksamt}[1]{\textbf{[#1~marks]}}
\newcommand{\markamt}[0]{\textbf{[1~mark]}}
\newcommand{\fillinMC}[1]{\fillinMCmath{\mbox{#1}}}
\newcommand{\fillinMCmath}[1]{\begin{tikzpicture}\draw circle [radius=0.5em];\end{tikzpicture}\ #1}
\newcommand{\fillinMCsoln}[1]{\fillinMCmathsoln{\mbox{#1}}}
\newcommand{\fillinMCmathsoln}[1]{\begin{tikzpicture}\draw[black, fill=blue] circle [radius=0.5em];\end{tikzpicture}\ #1}
\newcommand{\fillinMCAll}[1]{\fillinMCAllmath{\mbox{#1}}}
\newcommand{\fillinMCAllmath}[1]{\begin{tikzpicture}\draw (0,0) rectangle (1em,1em);\end{tikzpicture}\ #1}
\newcommand{\fillinMCAllsoln}[1]{\fillinMCAllmathsoln{\mbox{#1}}}
\newcommand{\fillinMCAllmathsoln}[1]{\begin{tikzpicture}\draw[black, fill=blue] (0,0) rectangle (1em,1em);\end{tikzpicture}\ #1}
\newcommand{\fillinblank}[1]{\fillinblankmath{\mbox{#1}}}
\newcommand{\fillinblankmath}[1]{\begingroup\setlength{\fboxsep}{1em}\setlength{\fboxrule}{2pt}\fbox{\LARGE\phantom{#1}}\endgroup}
\newcommand{\fillinblanksoln}[1]{\fillinblankmathsoln{\mbox{#1}}}
\newcommand{\fillinblankmathsoln}[1]{\begingroup\setlength{\fboxsep}{1em}\setlength{\fboxrule}{2pt}\fbox{{#1}}\endgroup}
\newcommand{\answerbox}[1][3\baselineskip]{%
\noindent\setlength{\fboxrule}{2pt}\framebox[0.7\linewidth]{%
\raisebox{2pt}[0pt][#1]{}%
}\par\medskip%
}
\usepackage{collcell,array}
\newcommand{\shiftdown}[1]{\smash{\raisebox{-.5\normalbaselineskip}{#1}}}
\newcolumntype{C}{>{\collectcell\shiftdown}c<{\endcollectcell}}
\newcommand{\shiftup}[1]{\smash{\raisebox{.5\normalbaselineskip}{#1}}}
%\fancyhead{}\fancyfoot[C]{\tiny This work is licensed under a \href{http://creativecommons.org/licenses/by/4.0/}{Creative Commons Attribution 4.0 International License. \ccby}\\ For license purposes, the author is the University of British Columbia.}\pagestyle{fancyplain}
\fancyhead{}\fancyfoot[C]{\tiny This work is licensed under a \href{http://creativecommons.org/licenses/by-nc/4.0/}{Creative Commons Attribution-NonCommercial 4.0 International License. \ccbync}\\ For license purposes, the author is the University of British Columbia.}\pagestyle{fancyplain}
\makeatletter
\def\@myproof#1[#2]{\trivlist \item [\hskip\labelsep{\bf #1 {#2}: }]}
\def\@begproof#1{\@ifnextchar[{\@myproof{#1}}{\@myproof{#1}[\unskip$\,$]}}
\newenvironment{solution}{\vskip -10pt\leavevmode\@begproof{Solution}}{\endtrivlist}
\makeatother
\date{}
\title{CPSC 320 2018W2: Assignment 2}
\hypersetup{
  pdfkeywords={},
  pdfsubject={},
  pdfcreator={Emacs 25.2.1 (Org mode 8.2.10)}}
\begin{document}

\maketitle
\vspace{-.5in}

Please follow the guidelines given in Assignment 1 for submission to
Gradescope, and for group collaboration. Remember to provide short
justifications for your answers. Submit by the deadline \textbf{Friday
February 1, 2019 at 10PM}.

If you are using the .tex file to prepare your answers, then for
questions where you need to select a circle you can
change \verb~\fillinMCmath~ to \verb~\fillinMCmathsoln~ for your choice of answer.

All of the problems in this assignment concern graphs. Throughout,
unless otherwise specified, graphs have $n$ nodes and $m$ edges.
A node's \emph{degree} is the number of edges incident on it. 
Assume that graphs have neither self-loops (edges from a node to itself), nor
multiple edges between the same two nodes.

A \emph{path} of length $k$ is a list of $k+1$ nodes, such that there is an edge
of the graph between each consecutive pair of nodes.  The \emph{shortest path}
between two nodes is a path with the minimum number
of edges. A \emph{simple path} repeats no node. A \emph{cycle} is a path that
starts and ends at the same node. A \emph{simple cycle} repeats no node
other than the starting/end node (which only appears twice).  

\medskip
\section{Graph Data Structures}
\label{sec-1}
Suppose that the edges of undirected graph $G = ([1..n],E)$ are
represented using \emph{adjacency lists}: for each $i$, there is a (not
necessarily sorted) linked list $A[i]$ of the nodes $j$ such that
$(i,j) \in E$.  Suppose furthermore that for each $i$, the degree of
$i$ is also stored, and so can be obtained in $\Theta(1)$ time.  With
these assumptions, what is the worst-case running time of the most
efficient algorithm you can think of to solve each of the following problems,
using basic graph algorithms familiar to you?

\begin{enumerate}
\item Given two nodes $i$ and $j$ of $G$, determine if they are adjacent.

$\begin{array}{lcl}
\fillinMCmathsoln{O(n)} & \qquad & \fillinMCmath{O(nm)} \\
\fillinMCmath{O(n\log n)} & \qquad & \fillinMCmath{O(n^2m)}\\
\fillinMCmath{O(n^2)} & \qquad & \fillinMCmath{O(nm^2)}\\
\fillinMCmath{O(n+m)} & \qquad & \fillinMC{None of these}
\end{array}$\\
\ans{
  Simply go through all the adjacency list of $i$ and see if there is the node $j$, this will 
  take at worst $O(n)$ time because the size of the adjacency list can be at most $n-1$.
}
\item Determine if $G$ is a star. An undirected graph is a \emph{star} if
for some node $i$ (called the star
centre), $E = \{(i,j) \;|\; i \neq j, 1\le j \le n\}$.  For example,
if $n=4$ then the graph with edges $\{(1,2), (2,3), 2,4)\}$ is a star
with centre 2. (Recall that for undirected graphs, the order of 
pair of nodes in an edge does not matter.)

$\begin{array}{lcl}
\fillinMCmathsoln{O(n)} & \qquad & \fillinMCmath{O(nm)} \\
\fillinMCmath{O(n\log n)} & \qquad & \fillinMCmath{O(n^2m)}\\
\fillinMCmath{O(n^2)} & \qquad & \fillinMCmath{O(nm^2)}\\
\fillinMCmath{O(n+m)} & \qquad & \fillinMC{None of these}
\end{array}$\\
\ans{
  Go through all nodes and check the degree of all, there must be only one node of degree $n-1$, all 
  the other nodes must have degree $1$. This will be $O(n)$ because we are jut looping through the nodes.
}

\item Determine if $G$ contains a triangle, that is, three
nodes $i,j,k$ such that $i$ and $j$ are adjacent,
$i$ and $k$ are adjacent, and $j$ and $k$ are adjacent.

$\begin{array}{lcl}
\fillinMCmath{O(n)} & \qquad & \fillinMCmath{O(nm)} \\
\fillinMCmath{O(n\log n)} & \qquad & \fillinMCmath{O(n^2m)}\\
\fillinMCmath{O(n^2)} & \qquad & \fillinMCmath{O(nm^2)}\\
\fillinMCmath{O(n+m)} & \qquad & \fillinMCsoln{None of these}
\end{array}$\\
\ans{
  This can be solve using a DFS algorithm with some modifications. We will go through each 
  node $i$ in $G$ (this will have a running time of $O(n)$). Then with each $i$ do a DFS 
  with a depth limit set to $3$. So if we do not find the node $i$ repeated at depth $3$ 
  then we know that there is no other nodes $j$ and $k$ that make a triangle 
  between $i$, $j$ and $k$. This DFS will have a worst running time of (n+m), so the 
  algorithm will have a complexity $O(n^2+nm)$ \\
}

\item Find a node $i$ of $G$ that minimizes $d(i) = \max_j \{
      \textit{dist}(i, j) \}$, where $\textit{dist}(i,j)$ is the
length of the shortest path connecting $i$ to $j$. (Again, it's
not obvious for this problem that there might not be a better
algorithm, but this seems to be the best possible with simple
graph algorithm techniques.)
\ans{
  To calculate $\textit{dist}(i,j)$ we will use BFS that will take complexity $O(n+m)$ 
  and there are only n choose 2 edges with different $(i,j)$ values, this is because 
  $(i,j)$ is the same that $(j,i)$. So wee need to run BFS n choose 2 edges, the complexity 
  would be $O(n^3+n^2m)$
}

$\begin{array}{lcl}
\fillinMCmath{O(n)} & \qquad & \fillinMCmath{O(nm)} \\
\fillinMCmath{O(n\log n)} & \qquad & \fillinMCmath{O(n^2m)}\\
\fillinMCmath{O(n^2)} & \qquad & \fillinMCmath{O(nm^2)}\\
\fillinMCmath{O(n+m)} & \qquad & \fillinMCsoln{None of these}
\end{array}$
\end{enumerate}

\section{Graph Properties}
\label{sec-2}
Each of the following problems presents a scenario about a graph and a
statement about that scenario. For each one, indicate by filling in
the appropriate circle whether:
\begin{itemize}
\item The statement is \textbf{ALWAYS} true, i.e., true in \emph{every} graph
matching the scenario.
\item The statement is \textbf{SOMETIMES} true, i.e., true in some graph
matching the scenario but false in another graph.
\item The statement is \textbf{NEVER} true, i.e., true in \emph{none} of the graphs
matching the scenario.
\end{itemize}


\begin{enumerate}
\item \textbf{Scenario:} An undirected graph containing two simple paths that
share no edges, each of length $k$ for some integer $k \geq
   2$. \textbf{Statement:} $n > k+1$.

$\fillinMCsoln{ALWAYS}$

$\fillinMC{SOMETIMES}$

$\fillinMC{NEVER}$\\
\ans{
  The smaller instance of the problem would be a graph with $5$ nodes all connected in a 
  line. The $k$ in this instance would be $2$ so the statement is true $5 > 2 + 1 $. 
  If we add one node to one size the $n$ will increase by $1$, the statement still true 
  $6 > 2 + 1 $. The only way to increase $k$ is by adding two nodes one on each end of 
  the two paths, this will increase the $k$ by $1$ and $n$ by two. Which preserves the 
  statement $7 > 3 + 1$.\\ \\
  There is no way that the $k$ is increased when the $n$ doesn't so by adding more nodes 
  two those two paths or to any part in the graph, the difference between $n$ and $k$ will 
  be larger every time.
}

\item \textbf{Scenario:} A rooted tree found by running DFS from some
node of a connected,
undirected graph with $n \geq 2$. \textbf{Statement:} Every
node in the tree has at most two children.

$\fillinMC{ALWAYS}$

$\fillinMCsoln{SOMETIMES}$

$\fillinMC{NEVER}$ \\
\ans{
  This depends on the topology of the graph and which metric the DFS uses to determine 
  the order in which adjacent nodes are visited. An example that proves that the statement
  is not always true, is with the definition of a star graph from the excercise \textbf{1.2}. \\ \\
  If we run DFS on a star graph $G$ starting from the center of the star, then our rooted tree 
  found by our DFS will have the root node with $n-1$ children. So if $G$ has $n <= 3$ the 
  statement will be true the DFS graph will have at most $2$ children. But if $n > 3$ then 
  the statement is false, i.e., for $n = 4$ the root node of the DFS graph will have $3$ 
  children.
}

\item \textbf{Scenario:} A weakly-connected but \textbf{not} strongly-connected,
directed graph with $n \geq 2$. \textbf{Statement:} A simple cycle of
length $n+1$ exists.

\textbf{SOLUTION}

$\fillinMC{ALWAYS}$

$\fillinMC{SOMETIMES}$

$\fillinMCsoln{NEVER}$\\
\ans{
  There can't be a simple cycle of length $n + 1$ in any graph, 
  because a cycle of length $n + 1$ needs to be connected with $n + 2$ nodes, so two nodes 
  will have to be repeated, on a simple cycle there can't be more repeated nodes rather 
  than the start/end node, so it is impossible that a \textbf{simple} cycle of length 
  $n+1$ exists in any directed graph.
}
\end{enumerate}

\section{Graph Application: Tutorial Assignments}
\label{sec-3}
We give definitions of two problems here; your task is to reduce the
first to the second. (After reading and understanding the two problem
definitions, and before continuing on to read more, make some notes
for yourself of what is needed to do a reduction and analyze its
running time, and then what is needed to show that the reduction is
correct. The rest of this problem breaks these parts down for you, but
it's good to get practice at breaking the parts down for yourself
too.)

\begin{itemize}
\item \textbf{Tutorial Assignment Problem (TAP)}. 
An instance of TAP is given by
\begin{itemize}
\item A list of tutorials $t_1, t_2, \ldots, t_k$ and their capacities
$c_1, c_2, \ldots, c_k$, such that $\sum_{j=1}^k c_j = n$.
\item A set $S_i$ of the tutorials that student $i$ can attend, where the total number of
students equals the total capacity $n$.
\end{itemize}
A \textbf{valid} solution is an assignment of students to tutorials such that 
\begin{itemize}
\item the total number of students assigned to tutorial $t_j$ is at most $c_j$, $1\le j \le k$,
\item each student $i$ is either unassigned or is assigned to a tutorial in $S_i$, $1 \le i \le n$, and
\item no student is assigned to more than one tutorial.
\end{itemize}
We can write a valid solution as a set $A$ of items
of the form $t_j$:$\{i_1,i_2,...\}$, where $i_1,i_2$, etc. are
the students assigned to tutorial $t_j$. Not all
students need be assigned to tutorials. A \textbf{good} solution is a
valid solution that maximizes the number of assigned students.

\item \textbf{Bipartite Matching Problem (BMP).} An instance of BMP is an undirected
bipartite graph $G=(X,Y,E)$ where $E \subseteq X \times Y$. That
is, the nodes of $G$ are partitioned into two sets $X$ and $Y$, and
all edges of $G$ connect a node of $X$ with a node of $Y$.  A \textbf{valid
solution} is a matching, i.e., a subset $M$ of $E$ in which each
node appears at most once. A \textbf{good} solution is a maximum matching.
\end{itemize}


\begin{enumerate}
\item Describe a reduction from TAP to BMP.

\textbf{Part (i):} Transform an instance $I$ of TAP to an instance $I'$ of BMP.\\
\ans{
  To transform the instance $I$ into $I'$ we need to create an undirected bipartite graph $G$. In this 
  graph, one partition will be filled with $n$ students nodes. The other partition will be $n$ nodes 
  that each one is the representation of a 'seat' in each tutorial, so for every tutorial $t_j$ 
  it will be $c_j$ nodes that we will denote as $seat_{t_jx}, 1<=x<=c_j$. In this way we ensure that there will be no more than $c_j$ students 
  assigned to the tutorial $t_j$. So the graph will have $n$ students and $n$ 'seats' of all the 
  tutorials. \\ \\
  The graph also need to have edges between students and 'seats', this edges will represent that a 
  student $i$ is capable of attend to a tutorial $t_j$. So to create this edges, we loop through the 
  set $S_i$ of each student $i$ and for every tutorial $t_j$, that $i$ can attend, we create an edge 
  between $i$ and every 'seat' node of the tutorial $t_j$ in $G$. \\
  This process will give us the graph $G$ that we can use to solve BMP.
}

\textbf{Part (ii):} Transform a solution $M'$ for instance $I'$ of BMP to
a solution $A$ for instance $I$ of TAP. \\
\ans{
  The solution $M'$ will be a set containing matches between students and 'seats' on tutorials. To 
  transform $M'$ in to the set $A$, that will contain a set of students for each tutorial $t_j$, we 
  need to loop through $M'$ and for each match $(i,seat_{t_jx})$ we will add $i$ to the set $t_j$ in 
  $A$. This process will give us a valid solution for the instance $I$ on the TAP problem.
}

\item Give the worst-case running time of part (i)
of your reduction, as a big-O function of $n$. \\
\ans{
  The worst-case running time, will be in the case that each student is capable of attending to all 
  the tutorials. This will give us a big-O $O(n^2)$ because for each student $i$ we need to add an edge 
  to all the tutorial seats which in the worst case will be $n$ seats for each $i$ student. 
}

\item Give the worst-case running time of part (ii)
of your reduction, as a big-O function of $n$.
\ans{
  We will need to loop through all of the matchings in the solution $M'$ and 
  for each one add each student to each corresponding tutorial, the length of $M'$ will be at worst 
  of length $n$. So the worst-case running time will be $O(n)$
}

\item There is a well-known algorithm for solving BMP with running time
$O(nm)$. (We won't cover the algorithm, but it can be
found in Chapter 7 of K\&T's textbook.) Combining this algorithm
with your reduction, what is worst-case running time of
your TAP algorithm, as a function of $n$? \\
\ans{
  On the worst case the instance generated instance $I'$ will have $2n$ nodes and $n^2$ edges, so if 
  the algorithm that computes the solution for BMP in $O(nm)$ then the runtime of computing $I'$ would 
  be $n^3$ since $m = n^2$, if we consider this in the reduction algorithm the whole running time 
  will be $n^2 + n^3 + n$, that will give us $O(n^3)$.
}

\item Explain why your reduction is correct.  Suppose that $M_{good}'$ is
a good solution (maximum matching) for $I'$ and let
$A_{good}$ be the transformed solution (assignment) obtained by part (ii) of
your reduction applied to $M_{good}$.

\begin{enumerate}
\item Show that $A_{good}$ is a valid solution for $I$. \\
\ans{
  The instance $I'$ is a completely valid instance for BMP, because we created an undirected 
  bipartite graph with the instance $I$ for the TAP. Therefore $M'_{good}$ must be a valid solution 
  for BMP. \\
  Then $M'_{good}$ must have a set of matchings between students and tutorial seats, in which each 
  student $i$ appears at most once, every tutorial seat $seat_{t_jx}$ is link to at most $1$ student, 
  and there are at most $c_j$ 'seats' of tutorial $t_j$. So by visiting each matching in the set 
  $M'_{good}$ we can create a set $A_{good}$ and for each student in $M'_{good}$ we add it to 
  the corresponding tutorial $t_j$ in $A_{good}$. Giving us at the end a valid solution with the 
  instance $I$ for TAP.
}

\item Show that $A_{good}$ is a good solution for $I$, i.e., maximizes
the number of assigned students. \\
\ans{
  If $M'_{good}$ is a good solution to BMP, meaning that it have the maximum matching in $I'$ which 
  represents students and tutorials seats, then $A_{good}$ will be a good solution for TAP, 
  because it will have the most number of students matched for a tutorial, while still being a 
  valid solution for TAP (proved on point 3.5.a).
}
\end{enumerate}
\end{enumerate}
\section{Graph Application: Network Connectivity}
\label{sec-4}
(Adapted from Problem 9, Chapter 3 of K\&T) Think of a communications
network  as a connected, undirected graph, where messages from one node $s$ to another node
$t$ are sent along paths from $s$ to $t$.  Nodes can sometimes fail.
If a node $v$ fails then no messages can be sent along edges incident
on $v$. A network is particularly vulnerable if failure of a
single node $v$ can cause two nodes, say $s$ and $t$, to become
disconnected. That is, when $v$ and its incident edges are removed
from the network, there is no path from $s$ to $t$.
We call such a node $v$ a \emph{vulnerability}.

\begin{enumerate}
\item Suppose that connected, undirected graph $G$ contains two nodes $s$ and $t$ such
that the shortest path between $s$ and $t$ is strictly greater than
$n/2$. Show that $G$ must contain a vulnerability. \\
\ans{
  In order to have a path between two nodes ($u$ and $v$) that is strictly greater than $n/2$ in $G$, 
  there must be an articulation node $w$ in the graph. We can see $w$ as a bottleneck where the 
  path must pass to reach $v$ from $u$. \\ \\
  Because $G$ is complete if there are no vulnerabilities all the paths from any node to 
  another node are at most $n/2$ in length. The only way to increase the length of the longest path 
  in $G$ is to add an extra node $x$ and connect it to a node which we called $w$, 
  this will increase the length of the longest path to be greater than $n/2$. 
  But this will introduce a vulnerability to $G$ in the node $w$.\\ \\
  So if $w$ is disconnected from the graph it will create at least two separate connected 
  components, in one of them it will be all the first part of the graph and in the other one it will 
  be the node $x$ that we connected to $w$. So we can't reach $x$ from any other node in the 
  $G$, hence $w$ will be a vulnerability.
}
\clearpage
\item Give an algorithm that, given as input $G$ and a node $s$ such that
there is some node $t$ for which the distance between $s$ and $t$
is guaranteed to be greater than $n/2$, finds a vulnerability. Your
algorithm should have running time $O(m + n)$. 
\begin{tabbing}
  xxxx\=xxxx\=xxxx\=xxxx\=xxxx\= \kill \color{ans}
  {\bf Algorithm} {\em Find-Vulnerability} ($V,E,s, depth=0$)\\ [0pt \color{ans}]
  \> // Initialize queue  \\
  \> let $n$ be size of $V$ \\
  \> let $q$ be a queue \\
  \> mark $s$ as visited \\
  \> set $s$.depth to $0$ \\
  \> $q$.enqueue($s$) \\
  \> // Iterate on the queue until is empty  \\
  \> while $q$ is not empty: \\
  \>\> let $t$ be $q$.dequeue \\
  \>\> // Go through all neighbours of $t$  \\
  \>\> For all neighbours $v$ of $t$: \\
  \>\>\> if $v$ is not marked: \\
  \>\>\>\> // If the path from $s$ to $v$ has a length greater than $n/2$ return $t$ \\
  \>\>\>\> if $t$.depth + 1 $ > n/2$: \\
  \>\>\>\>\> return $t$ \\
  \>\>\>\> // If not add $v$ to queue  \\
  \>\>\>\> mark $v$ as visited \\
  \>\>\>\> set $v$.depth to $t$.depth + 1 \\
  \>\>\>\> $q$.enqueue($v$) \\
  \> // No vulnerability found \\
  \> return null \\
\end{tabbing}
\ans{
  The algorithm above has a running time $O(m+n)$ because it simply does a BFS from the node 
  $s$, keeping track of the depth and when the depth reaches a level greater than $n/2$ 
  then the node that creates that path in that depth, must be a vulnerability. 
}
\end{enumerate}
% Emacs 25.2.1 (Org mode 8.2.10)
\end{document}
